\Askhsh{28314_SOLUTION}{1}
\begin{ylh}{1}
    \item [Συνέχεια, Παραγωγισιμότητα, Μονοτονία, Σύνολο τιμών συνάρτησης]
\end{ylh}
ΛΥΣΗ
\begin{alist}
    \item Το πεδίο ορισμού της $f$ είναι το $A = [1, +\infty)$.
    Επειδή η συνάρτηση $f$ είναι συνεχής στο πεδίο ορισμού της, θα είναι συνεχής και στο $x_0 = 1$
    οπότε θα ισχύει $\lim_{x \to 1^+} f(x) = f(1)$ ή $\lim_{x \to 1^+} e^{\frac{1}{\lambda x + 1}} = 0$ (1).
    Αν $\lambda \neq -1$ τότε $\lambda + 1 \neq 0$ και $\lim_{x \to 1^+} e^{\frac{1}{\lambda x + 1}} = e^{\frac{1}{\lambda+1}} > 0$, που είναι άτοπο λόγω της (1).
    Αν $\lambda = -1$ τότε $\lim_{x \to 1^+} (\lambda x + 1) = \lim_{x \to 1^+} (1 - x) = 0$ και $1 - x < 0$ για κάθε $x > 1$, άρα
    \[ \lim_{x \to 1^+} \frac{1}{\lambda x + 1} = \lim_{x \to 1^+} \frac{1}{1 - x} = -\infty \]
    οπότε
    \[ \lim_{x \to 1^+} e^{\frac{1}{\lambda x + 1}} = \lim_{x \to 1^+} e^{\frac{1}{1 - x}} = 0. \]
    Επομένως $\lambda = -1$. \monades{X}

    \item Για $\lambda = -1$ έχουμε
    \[ f(x) = \begin{cases} e^{\frac{1}{1-x}}, & x > 1 \\ 0, & x = 1 \end{cases} \]
    Για κάθε $x > 1$ η $f$ είναι παραγωγίσιμη με
    \[ f'(x) = \left(e^{\frac{1}{1-x}}\right)' = e^{\frac{1}{1-x}} \cdot \left(\frac{1}{1-x}\right)' = e^{\frac{1}{1-x}} \cdot \frac{1}{(1-x)^2}. \]
    Στο $x_0 = 1$ έχουμε:
    \[ \lim_{x \to 1^+} \frac{f(x)-f(1)}{x-1} = \lim_{x \to 1^+} \frac{e^{\frac{1}{1-x}}-0}{x-1} \text{ (2)}. \]
    Θέτουμε $\frac{1}{1-x} = u$ οπότε $\lim_{x \to 1^+} u = \lim_{x \to 1^+} \frac{1}{1-x} = -\infty$ και το όριο (2) γίνεται
    \[ \lim_{u \to -\infty} (-ue^u) = \lim_{u \to -\infty} \frac{-u}{e^{-u}} \xrightarrow{\left(\frac{+\infty}{+\infty}\right) \text{ de l' Hospital}} \lim_{u \to -\infty} \frac{(-u)'}{(e^{-u})'} = \lim_{u \to -\infty} \frac{-1}{-e^{-u}} = \lim_{u \to -\infty} \frac{1}{e^{-u}} = \lim_{u \to -\infty} e^u = 0. \]
    Επομένως η συνάρτηση $f$ είναι παραγωγίσιμη στο $x_0 = 1$ με $f'(1) = 0$.
    Επομένως
    \[ f'(x) = \begin{cases} e^{\frac{1}{1-x}} \cdot \frac{1}{(1-x)^2}, & x > 1 \\ 0, & x = 1 \end{cases} \] \monades{X}

    \item Λόγω του ερωτήματος (β), για κάθε $x > 1$ είναι $f'(x) = e^{\frac{1}{1-x}} \cdot \frac{1}{(1-x)^2} > 0$ και επειδή η συνάρτηση $f$ είναι συνεχής στο $x_0 = 1$, η $f$ θα είναι γνησίως αύξουσα στο $A = [1, +\infty)$.
    Αφού η $f$ είναι γνησίως αύξουσα στο $A = [1, +\infty)$ θα παρουσιάσει στο $x_0 = 1$ ελάχιστο το $f(1) = 0$. \monades{X}

    \item Είναι $\lim_{x \to +\infty} f(x) = \lim_{x \to +\infty} e^{\frac{1}{1-x}} = 1$, αφού $\lim_{x \to +\infty} \frac{1}{1-x} = 0$. Επειδή η $f$ είναι γνησίως αύξουσα και συνεχής στο $A = [1, +\infty)$, το σύνολο τιμών της θα είναι το
    \[ f(A) = [f(1), \lim_{x \to +\infty} f(x)) = [0, 1). \] \monades{X}
\end{alist}